\newcommand{\figuraAceleracionCentrifugaRecipiente}{%
\begin{figure}[H]
    \centering
    \begin{tikzpicture}[>=Latex,font=\small,line cap=round,line join=round]
        \node[font=\bfseries] at (7.0,6.25)
            {\texthalo{Interpretación física en el recipiente rotante}};

        \draw[very thick] (2.25,5.45)--(2.25,0.65)--(11.75,0.65)--(11.75,5.45);
        \fill[blue!12] (2.33,0.73) rectangle (11.67,4.35);
        \draw[very thick,blue!65!black] (2.25,4.35)--(11.75,4.35);

        \draw[densely dashed,black!65] (7.0,0.65)--(7.0,5.75);
        \node[anchor=south east] at (6.88,5.42)
            {\texthalo{eje de rotación}};

        \fill[orange!80!black] (9.45,3.05) circle (0.14);
        \node[anchor=south] at (9.45,3.55)
            {\texthalo{elemento de fluido}};
        \draw[->,thick] (7.0,3.05)--(9.45,3.05)
            node[midway,below] {\texthalo{$\vec r=r\hat e_r$}};
        \draw[->,ultra thick,red!75!black] (9.45,3.05)--(11.25,3.05)
            node[midway,above] {\texthalo{$\vec a_{\mathrm{cf}}$}};
        \node[red!75!black,anchor=west] at (9.55,2.15)
            {\texthalo{$|\vec a_{\mathrm{cf}}|=\Omega^2r$}};

        \draw[->,ultra thick,purple!75!black] (7.0,1.55)
            arc[start angle=90,end angle=-235,radius=0.58];
        \node[purple!75!black,anchor=west] at (7.78,1.55)
            {\texthalo{$\vec\Omega$}};

        \draw[->,very thick,red!75!black] (3.25,2.15)--(4.05,2.15);
        \draw[->,very thick,red!75!black] (4.55,2.15)--(5.75,2.15);
        \node[align=center] at (4.50,1.45)
            {\texthalo{la magnitud aumenta}\\
             \texthalo{linealmente con $r$}};

        \fill[orange!80!black] (7.0,4.35) circle (0.10);
        \node[anchor=south west] at (7.10,4.43)
            {\texthalo{$r=0\Rightarrow\vec a_{\mathrm{cf}}=\vec 0$}};
        \node[align=center,text=black!65] at (7.0,0.12)
            {\texthalo{La aceleración es radial y apunta siempre hacia afuera del eje.}};
    \end{tikzpicture}
    \caption{Aceleración centrífuga dentro de un recipiente que gira con
    velocidad angular constante. Se anula sobre el eje de rotación y crece
    linealmente con la distancia radial: su magnitud es $\Omega^2r$ y su
    dirección es $\hat e_r$.}
    \label{fig:aceleracion-centrifuga-recipiente}
\end{figure}%
}